% =============================================================================
%  Sección 2: Crisis en las Ciencias
% =============================================================================
\section{Crisis en las Ciencias}

\subsection{Conflicto entre Ciencias}

No son pocos los debates que existen entre la división de Ciencias Sociales y Ciencias exactas 
o Ciencias Matemáticas. Dicha comparación ha ido evolucionando y escalado desde principios del Siglo XX.
Tomaremos como ejemplo el caso de la elevación a Ciencia de la Sociología, ya que anteriormente esta 
era considerada una rama de la Filosofía. Karl Popper \cite{popper_miseria_2014} hace una comparación interesante; 
¿Dónde radica su carácter como ciencia? y más importante aún ¿Por qué la comparación entre una ciencia de
carácter meramente matemático con una social? Ya desde principios del siglo XIX Augusto
Comte fue el máximo exponente de elevar a la observación de las sociedades como una
ciencia, aplicando el método positivista (es decir el método científico) para explicar, estructurar
y proponer dos ramas: La estática social y la dinámica social. Sin embargo, no es hasta
principios del siglo XX con Emilie Durkheim \cite{durkheim_reglas} quién dotó a
la ciencia de una fama y aplicación internacional. Dándole un carácter más rígido, estructurado
y un objetivo propio. \\

Sin embargo, el ser humano, en su afán de querer controlar todo lo que le rodea, creó y sistematizó una serie de pasos 
de manera exacta, predecible y replicable basándose en un sistema numérico casi universal: \termino{Las Matemáticas}. 
A través de ello, se inventaron diferentes artefactos de medición como lo son calendarios, relojes, escuadras y más recientemente 
sistemas de cómputo. Pero estos se quedan en meras herramientas de observación, medición y cálculo, no en un control del factor 
físico del espacio y tiempo. Una de las carencias más sonadas en el siglo XXI, es que el humano ha ido relegando poco a poco 
el razocinio análisis y entendimiento a sistemas de cómputo al pasar de los años. Este problema aunque se ha popularizado en los 
últimos años, no es más que un reflejo de la división y divorcio que ha existido desde siempre entre lo social y lo exacto. 
Tomemos como referencia la siguiente cita:
\begin{nopagebreak}
\citaclave{
    \textit{``A computer can never be held accountable, \\ therefore a computer must never make a management decision.''}~\citep{ibm_training_1979} 
    \\
    {\color{gray}\hrule height 0.4pt} % Línea divisoria tenue
    \textit{ \\ 
        ``Una computadora nunca puede ser responsabilizada, por lo tanto, una computadora nunca debe tomar una decisión de gestión.'' (Traducción al español)}
    }
\end{nopagebreak}

Tomado de un manual de cómputo básico para sistemas electrónicos y digitales, el peso es contundente; 
jerarquizar las ciencias es un error sistemático en el análisis de resultados que requieren una intervención multidisciplinario. 
\cite{lefebvre_rural} Este escenario tiende a ser el más común en la mayoría de los problemas a escala Macro, 
si entendemos que el avance y correcta implementación de la Ciencia y la tecnología es un proceso que va de la mano. 
Este problema resuena más en los últimos años, pues la Inteligencia Artificial o \termino{IA} ha tomado una fuerza descomunal en 
el reciente lustro (2021-2026) y expone una fractura del sistema que lo originó: El sistema capitalista no busca la unión o la 
democratización de procesos equitativos y cualitativos; busca la eficiencia en los procesos al menor costo posible \cite{bauman_globalizacion_1999}.
Lo que se traduce en la acción de las empresas por minimizar costos operativos y mano de obra humana por el proceso sistemático de 
una fría computadora; chocando directamente con la cita de IBM, donde no se le puede culpar a una computadora por la toma de decisiones.

\subsection{Urbanismo como afectado}
Este divorcio de las ciencias, responde más a un cambio de paradigma de quienes tienen el poder de decisión en el sistema académico, político y 
económico; es decir la capacidad de explorar/investigar, financiar y aplicar. En estudios muy especializados como lo pueden ser la creación de una nueva 
tarjeta gráfica (a nivel computacional) o el estudio cualitativa de un pueblo originario, no hay mayor afectación, ya que pocas veces se convierten 
en estudios multidisciplinarios. Pero en el caso del Urbanismo como disciplina, donde se congenian materias como la Arquitectura, la Sociología, Psicología y 
más recientemente la Computación y el análisis de datos mediante software especializado como lo son los \termino{Sistemas de Información Geográfica (SIG)}, 
se convierte en la analogía de la \termino{Bola de Nieve:} Mientras más rueda, más grande se va haciendo; y en algún punto golpeará a un objeto, dejando ver 
las fisuras que no fueron atendidas en su momento. El mismo Durkheim menciona que no tratar a los hechos sociales como cosas medibles, es un error sistemático. 
Y las fisuras o problemas como los on la violencia, delincuencia y el desorden, son los indicadores que necesita el analista para identificar los problemas 
de su sistema sociológico \cite{durkheim_reglas}. \\
En el caso singular de el transporte, se ha cometido muchas veces la falta de este enfoque multidisciplinario. Puesto que no es solo el movimiento de personas de 
Punto A - Punto B, involucra otros factores clave como lo son la seguridad de sus vialidades, los espacios de convivencia, la eficiencia de una red bien conectada, 
la robustez de esta red ante fallas ocasionadas por el medio ambiente (terremotos, inundaciones, etc.), fallas en el entorno social (disturbios, manifestaciones civiles) o 
inconveniencias en el mismo sistema (fallos económicos, de logística o infraestructura). Tal es el caso de los informes más recientes \cite{ramirez2025analisis}, donde se rescata la intervención 
en temas de sostenibilidad y desarrollo ecológico, pero se ignora la inmensidad de la demanda o del sistema geográfico al cual se sirve, al igual que las necesidades particulares 
de la población que hará uso de ellos. \textbf{De nada nos sirve un sistema ecológico y sostenible si es ineficiente en términos de robustez ante fallas, capacidad operativa y logística geográfica. 
Lamentablemente, esto es muy común en una urbe como la \zmvm.}